\begin{quote}\small
\noindent\textbf{Executive summary.} This paper benchmarks 23 emerging economies over 1990--2024 across three lenses: causal macro identification, sovereign credit pricing, and tradable emerging-market factor construction. The headline composite KPI score aggregates 39 development indicators across six sectors and places Singapore first ($+2.93\,\sigma$), China second ($+2.33\,\sigma$) and Niger last ($-2.58\,\sigma$); China carries the largest absolute improvement in the panel ($+2.36\,\sigma$ since 1990); Russia is the only panel member whose composite \emph{declines} over the period ($-0.29\,\sigma$). \emph{Causal identification} via Jord\`a local projections is noisy at this $N \times T$ sample; the Dumitrescu--Hurlin panel Granger test recovers significant effects of energy on poverty ($p = 0.027$) and health on Gini ($p = 0.046$); the Westerlund test confirms cointegration between the composite and log GDP per capita at $p = 0.009$. \emph{Sovereign credit} is the empirical headline: a pooled ordered probit of S\&P ratings on the six sector scores delivers out-of-sample MAE $= 1.84$ notches and AUC $= 0.81$ for the investment-grade binary --- $8$ percentage points above the macro-only baseline. A Bayesian hierarchical ordered-logistic counterpart localises the signal to energy ($\mu_\beta = +0.92$) and health ($\mu_\beta = -0.92$); the four other sectors have credible intervals that cover zero. A Lasso of the rating on sector scores selects R\&I, energy, education, housing-living and security-stability at the CV-optimal penalty; the equal-weighted and Lasso-weighted composites have Kendall $\tau = 0.91$ at the latest-year cross-section. A Cox proportional-hazards analysis on time-to-investment-grade gives an education hazard ratio of $1.49$ ($p = 0.16$) --- directionally positive but underpowered at $n = 21$. \emph{Tradable performance} fails. A KPI-sorted long-only top-quintile EM equity portfolio earns Sharpe $0.15$ over 2007--2024 against the MSCI EM benchmark's $0.12$, sitting at the $0^{\text{th}}$ percentile of the 5\,000 random-portfolio distribution (mean random Sharpe $0.28$). Extending the test to credit instruments does not rescue the signal: the cross-sectional KPI improvement does not predict EMB returns ($t = 0.73$, $R^2 = 0.003$, $n = 204$ months), and a long-upgrade / short-downgrade overlay loses $9.3\%$ annualised. The clean reading: KPIs are slow-moving and price into the cross-sectional level of sovereign ratings, not into monthly equity or credit returns. The paper is structured to make that point.
\end{quote}

\vspace{0.6em}\hrule\vspace{1.0em}
